%% This is file sample-sigconf.tex,
%% generated with the docstrip utility.
%%
%% The original source files were:
%%
%% samples.dtx  (with options: `sigconf')
%% 
%% IMPORTANT NOTICE:
%% 
%% For the copyright see the source file.
%% 
%% Any modified versions of this file must be renamed
%% with new filenames distinct from sample-sigconf.tex.
%% 
%% For distribution of the original source see the terms
%% for copying and modification in the file samples.dtx.
%% 
%% This generated file may be distributed as long as the
%% original source files, as listed above, are part of the
%% same distribution. (The sources need not necessarily be
%% in the same archive or directory.)
%%
%% The first command in your LaTeX source must be the \documentclass command.
\documentclass[sigconf,a4paper]{acmart}

%%
%% \BibTeX command to typeset BibTeX logo in the docs
\AtBeginDocument{%
  \providecommand\BibTeX{{%
    \normalfont B\kern-0.5em{\scshape i\kern-0.25em b}\kern-0.8em\TeX}}}

%% Rights management information.  This information is sent to you
%% when you complete the rights form.  These commands have SAMPLE
%% values in them; it is your responsibility as an author to replace
%% the commands and values with those provided to you when you
%% complete the rights form.
\setcopyright{acmcopyright}
\copyrightyear{2022}
\acmYear{2022}
%\acmDOI{10.1145/1122445.1122456}

%% These commands are for a PROCEEDINGS abstract or paper.
\acmConference[DistAlg25/26]{The first project delivery of DistAlg25/26}{2026}{Faculdade de Ciências e Tecnologia, NOVA University of Lisbon, Portugal}
\acmBooktitle{The Projects of DistAlg - first project, 2026, Faculdade de Ciências e Tecnologia, NOVA University of Lisbon, Portugal}
%\acmPrice{15.00}
%\acmISBN{978-1-4503-XXXX-X/18/06}


%%
%% Submission ID.
%% Use this when submitting an article to a sponsored event. You'll
%% receive a unique submission ID from the organizers
%% of the event, and this ID should be used as the parameter to this command.
%%\acmSubmissionID{123-A56-BU3}

%%
%% The majority of ACM publications use numbered citations and
%% references.  The command \citestyle{authoryear} switches to the
%% "author year" style.
%%
%% If you are preparing content for an event
%% sponsored by ACM SIGGRAPH, you must use the "author year" style of
%% citations and references.
%% Uncommenting
%% the next command will enable that style.
%%\citestyle{acmauthoryear}

%%
%% end of the preamble, start of the body of the document source.
\begin{document}

%%
%% The "title" command has an optional parameter,
%% allowing the author to define a "short title" to be used in page headers.
\title{Evaluating Membership and Broadcast Protocol Combinations in Gossip-Based Systems}

%%
%% The "author" command and its associated commands are used to define
%% the authors and their affiliations.
%% Of note is the shared affiliation of the first two authors, and the
%% "authornote" and "authornotemark" commands
%% used to denote shared contribution to the research.
\author{Miguel Mestre}
\authornote{Student number 73018}
\email{mfg.mestre@campus.fct.unl.pt}
\affiliation{%
  \institution{MEI, DI, FCT, UNL}
}

\author{Alexandre Santos}
\authornote{Student number 72970}
\email{aff.santos@campus.fct.unl.pt}
\affiliation{%
  \institution{MEI, DI, FCT, UNL}
}

\author{Nelson}
\authornote{Student number 11111. Marbles did not understood the assignment but he did use an LLM to make a song related with the project.}
\email{marbles@cool.se}
\affiliation{%
  \institution{MEI, DI, FCT, UNL}
}

%%
%% By default, the full list of authors will be used in the page
%% headers. Often, this list is too long, and will overlap
%% other information printed in the page headers. This command allows
%% the author to define a more concise list
%% of authors' names for this purpose.
\renewcommand{\shortauthors}{Miguel, Alexandre, and Nelson.}

%%
%% The abstract is a short summary of the work to be presented in the
%% article.
\begin{abstract}
  Gossip-based protocols have emerged as an effective approach to implement scalable and resilient broadcast mechanisms in distributed systems. In such systems, nodes rely on membership protocols to maintain partial views of the network, which directly influence the efficiency and reliability of message dissemination. The performance of gossip-based broadcast depends on the interaction between membership and broadcast strategies, particularly under varying system conditions.

Several approaches can be used to construct and maintain these overlays, each offering different trade-offs in terms of reliability, latency, and communication overhead. In this paper, we implement and experimentally evaluate combinations of membership protocols, namely Cyclon and HyParView, with broadcast protocols such as Flood and Eager-Push Gossip, as well as a custom-designed gossip-based solution. Our results show that different protocol combinations exhibit distinct performance characteristics, highlighting trade-offs between reliability and latency, and demonstrating that the choice of membership protocol plays a key role in determining the effectiveness of broadcast mechanisms.
\end{abstract}


%%
%% This command processes the author and affiliation and title
%% information and builds the first part of the formatted document.
\maketitle

\section{Introduction}

Modern distributed systems often need to send a message to all nodes
participating in the system, an operation known as broadcast. The simplest
way to do this is flooding, where every node that receives a message
immediately forwards it to all of its neighbors. This guarantees that every
node will eventually receive the message, but it comes at a steep cost.
Because every node forwards to every neighbor, the total number of messages
in the network grows very quickly as the system scales, filling up network
links and forcing each node to spend time processing and discarding large
numbers of duplicate messages. This extra work does not go unnoticed: the
congestion and processing overhead introduced by flooding ultimately manifests
as increased end-to-end latency, meaning messages take longer to arrive not
because the network is slow, but because it is overloaded.

A more efficient alternative is gossip-based broadcast, where instead of
forwarding to every neighbor, each node picks only a small random subset of
neighbors to forward to. The size of this subset is controlled by a parameter
called the fanout. At first glance it might seem that reducing the number of
forwarding targets would significantly hurt reliability, but in practice,
if the fanout is well calibrated, the message still reaches the vast majority
of nodes in the system. The core argument is therefore that the overhead
flooding introduces does not meaningfully improve reliability over a well
tuned gossip protocol, it only adds congestion and latency.

For gossip to work well, nodes need a way to know which other nodes exist in
the system and to maintain a set of neighbors to gossip with. This is the
role of an overlay network protocol. In this work we consider two such
protocols: Cyclon and HyParView. Cyclon maintains a single list of neighbors
and periodically shuffles it with other nodes to keep the view random and
fresh. It is simple and works well in stable conditions. HyParView takes a
different approach and maintains two lists: an active view of neighbors that
are currently in use, and a passive view of backup neighbors that are kept
in reserve. When an active neighbor leaves or stops responding, HyParView
can immediately replace it with a node from the passive view without waiting
for the next maintenance cycle. This makes HyParView more resilient when
nodes frequently join and leave the network, a condition known as churn.
In scenarios with little churn however, the simpler design of Cyclon is
expected to introduce slightly less overhead and therefore slightly lower
latency.

On top of these overlay protocols we implement a hybrid gossip broadcast
protocol. Rather than always forwarding the full message to neighbors, our
protocol splits neighbors into two groups per message. A small subset called
eager targets receives the full message right away. Another subset called
lazy targets receives only a small announcement indicating that the sender
has a message available. If a lazy target has not yet seen that message, it
can request the full content, effectively pulling it on demand. This reduces
unnecessary full message transmissions when nodes have already received the
message through another path. Additionally, our protocol monitors how many
duplicate messages it is receiving and uses this information to dynamically
adjust the eager fanout over time, growing it when too few nodes seem to be
receiving messages and shrinking it when too many duplicates are observed.
We expect this adaptive behavior to reduce message overhead compared to a
static eager push approach, though the dynamic adjustment of the fanout may
occasionally cause short periods of reduced delivery reliability.

In this paper we experimentally study these three aspects: whether flooding
truly introduces more latency than gossip for similar reliability levels,
whether HyParView outperforms Cyclon under high churn while Cyclon holds
an edge in stable networks, and whether our hybrid adaptive protocol
improves on static gossip in terms of overhead while remaining competitive
in reliability. The remainder of this paper is organized as follows.
Section~\ref{sec:related} discusses related work on overlay networks and
gossip protocols. Section~\ref{sec:implementation} describes our
implementation in detail. Section~\ref{sec:evaluation} presents and
discusses our experimental results. Finally,
Section~\ref{sec:conclusions} concludes the paper.

\section{Related Work}
\label{sec:related}

This section reviews prior work that is necessary to contextualize the
protocol combinations studied in this paper. We first discuss random
overlay construction protocols, then gossip-based broadcast protocols,
and finally prior experimental studies that compare dissemination
strategies in similar settings.

\subsection{Random Overlay Networks}

Random overlay networks are a common building block for scalable
distributed systems because they provide each node with only a partial
view of the network while still preserving good connectivity and
short communication paths. Instead of requiring every node to know all
other participants, these protocols maintain a small and continuously
refreshed neighbor set, which is enough to support robust dissemination
and membership management at scale.

Cyclon is one of the best-known peer-sampling protocols in this space.
Its main goal is to maintain a continuously evolving random overlay by
periodically exchanging small neighbor samples between nodes. This
periodic shuffling process keeps views fresh, removes stale entries,
and helps preserve overlay randomness over time. Because of its simple
design, Cyclon is often used as a baseline membership mechanism in
experimental studies of gossip-based systems.

HyParView was proposed to improve robustness under churn while still
maintaining a lightweight overlay. In contrast with protocols that rely
on a single partial view, HyParView separates neighbor management into
two structures: a small active view used for communication and a larger
passive view used as backup. This distinction allows the protocol to
react quickly to failures by replacing lost active neighbors from the
passive view, while periodic maintenance keeps the passive view updated.
As a result, HyParView is generally regarded as better suited to highly
dynamic environments.

These two protocols represent different design choices in overlay
construction. Cyclon emphasizes simplicity and randomness through
periodic exchanges, whereas HyParView combines reactive and periodic
mechanisms to improve failure resilience. This makes them natural
candidates for comparison when studying how the choice of membership
layer affects broadcast performance.

\subsection{Gossip-Based Broadcast}

Gossip-based broadcast protocols were introduced as a scalable
alternative to deterministic flooding. In flooding, each node forwards a
received message to all its neighbors. This strategy is simple and
achieves full dissemination in connected overlays, but it generates a
large number of duplicate transmissions and places unnecessary load on
the network. As the system grows, this redundancy increases overhead and
may also increase end-to-end latency due to congestion and repeated
processing.

Eager-push gossip reduces this cost by forwarding each message only to a
subset of neighbors, usually chosen at random. The size of this subset,
commonly called the fanout, determines the trade-off between reliability
and overhead. When the fanout is too small, messages may fail to reach
the whole system. When it is too large, the protocol starts behaving
closer to flooding and produces many duplicates. A well-chosen fanout
therefore allows gossip-based broadcast to retain high coverage while
significantly reducing communication costs.

Several works also explore hybrid dissemination strategies. Instead of
always sending full payloads, some protocols first advertise the
existence of a message and only transmit the full content when needed.
This idea appears in lazy-push and pull-based designs, where message
identifiers are disseminated separately from message bodies. Such
approaches can reduce redundant full-message transmissions, especially
when many nodes have already received the same content through different
paths.

Plumtree is one of the best-known protocols combining eager and lazy
strategies. Its design attempts to exploit the low latency of tree-based
dissemination while preserving the resilience of gossip through a backup
mechanism. This line of work shows that combining dissemination modes
can provide better trade-offs than purely eager or purely lazy schemes.

Our hybrid protocol is inspired by this family of approaches. However,
instead of fixing the balance between eager and lazy dissemination once
and for all, we adapt the eager fanout dynamically according to the
observed duplication rate. This design choice aims to let the protocol
react to the conditions of the overlay rather than relying on a single
static parameter.

\subsection{Experimental Studies of Gossip Protocols}

A substantial body of prior work evaluates gossip-based dissemination
protocols experimentally, typically focusing on metrics such as
reliability, dissemination latency, message overhead, and robustness
under churn. These studies repeatedly show that dissemination behavior
cannot be analyzed independently from overlay construction. The quality
of the membership layer influences path diversity, redundancy, and the
ability to recover from failures, all of which affect broadcast
performance.

Prior experimental comparisons also show that flooding tends to achieve
high reliability at the cost of excessive overhead, whereas probabilistic
gossip often achieves comparable coverage with fewer transmissions when
its parameters are well tuned. Similarly, overlays designed for dynamic
environments tend to outperform simpler overlays under churn, even if
they may introduce some additional maintenance cost in stable scenarios.

The goal of our work is aligned with this experimental tradition.
Rather than proposing a purely theoretical argument, we implement
different combinations of membership and broadcast protocols and analyze
their behavior under common dissemination metrics. In this sense, our
study is not only about the protocols in isolation, but also about the
interaction between overlay maintenance and message dissemination.

\section{Implementation}
\label{sec:implementation}

This section describes the implementation of the system developed for
this project. The goal of the implementation was to make it possible to
experimentally compare multiple combinations of membership and broadcast
protocols within a common execution framework. To achieve this, we
organized the system into two main layers: a membership layer, which is
responsible for maintaining the overlay network, and a broadcast layer,
which uses that overlay to disseminate messages. We first present the
overall organization of the implementation, then describe the membership
protocols, and finally discuss the broadcast protocols and the custom
hybrid solution.

\subsection{System Organization}

The system was implemented using the Babel framework, which supports the
composition of protocols through message handlers, timers, channel
events, and inter-protocol notifications. This model was particularly
useful for the project because it allowed us to separate membership
logic from broadcast logic while still enabling communication between
both layers.

The membership protocols expose changes in the local neighborhood through
notifications that inform the broadcast layer when neighbors become
available or unavailable. This design keeps the broadcast protocols
independent from the internal details of overlay maintenance and makes
it possible to combine the same broadcast algorithm with different
membership mechanisms.

At a high level, every node executes one membership protocol and one
broadcast protocol. The membership protocol builds and maintains the
local partial view of neighbors. The broadcast protocol relies on that
view to forward messages. This separation mirrors the conceptual
structure of the project and simplifies both experimentation and code
reuse.

\subsection{Membership Layer}

\subsubsection{Cyclon}

Our Cyclon implementation follows the standard peer-sampling approach
based on partial views and periodic shuffling. Each node maintains a
view containing neighbor entries associated with an age. At regular
intervals, the node increments the age of its entries, selects one
neighbor, and exchanges a sample of entries with that node. This process
gradually refreshes the view and helps preserve randomness in the
overlay.

The main advantage of Cyclon in our implementation is its simplicity.
Since all view maintenance is driven by periodic exchanges, the protocol
is compact and easy to integrate with the rest of the system. However,
this same design means that recovery from failures is less immediate,
because failed nodes are only removed or replaced when subsequent
maintenance actions occur.

\subsubsection{HyParView}

Our HyParView implementation extends the simpler peer-sampling approach
by maintaining two views: an active view and a passive view. The active
view contains the neighbors currently used for communication, while the
passive view contains backup nodes that can be promoted when necessary.

When a node joins the system, it contacts a known participant and sends a
join request. The contacted node inserts the newcomer into its active
view and propagates a \texttt{FORWARDJOIN} message through the overlay
using a random walk controlled by a time-to-live parameter. Intermediate
nodes may insert the new node into their passive view, while the final
node in the walk inserts it into its active view. This mechanism helps
distribute knowledge of the joining node across the network while
keeping active views small.

When an active connection fails, the protocol immediately tries to
replace the failed node using an entry from the passive view. This is
one of the central design goals of HyParView and the main reason it is
expected to perform better under churn. The passive view itself is
refreshed periodically through shuffle messages that exchange neighbor
information across the overlay.

\subsubsection{Why Compare Cyclon and HyParView}

These two membership protocols were implemented because they represent
different philosophies of overlay maintenance. Cyclon relies on a single
randomized structure maintained through periodic exchanges, while
HyParView relies on a small communication view plus a larger backup
view, combining reactive and periodic mechanisms. By implementing both,
we can observe how broadcast protocols behave when the overlay is either
simpler and lighter or more explicitly designed for fault recovery.

\subsection{Broadcast Layer}

\subsubsection{Flood Broadcast}

The flood broadcast implementation is the simplest dissemination
mechanism in the system. When a node first receives a message, it
delivers it locally and forwards it to all current neighbors except,
when applicable, the sender from which the message was received. To
avoid infinite retransmission loops, each node stores the identifiers of
messages it has already processed and discards duplicates.

This protocol serves as a baseline. Its expected strength is very high
reliability in connected overlays. Its expected weakness is the large
number of duplicate transmissions generated by forwarding to all
neighbors.

\subsubsection{Eager-Push Gossip}

In eager-push gossip, a node forwards each newly received message only
to a randomly selected subset of its neighbors. The size of this subset
is the fanout. As in flood broadcast, duplicate messages are ignored
through the use of message identifiers.

Compared with flooding, this protocol introduces an explicit trade-off.
Forwarding to fewer neighbors reduces overhead, but also makes
dissemination depend more strongly on the quality of the overlay and on
the selected fanout. This protocol is therefore useful to study the cost
of probabilistic dissemination under different membership layers.

\subsubsection{Hybrid Adaptive Gossip}

The custom protocol developed for this project combines eager and lazy
dissemination. When a node first receives a message, it divides the
available neighbors into two groups. A subset of neighbors receives the
full message immediately and forms the eager set. The remaining selected
neighbors receive only a lightweight notification containing the message
identifier and form the lazy set.

If a node receives a lazy notification for a message it does not yet
have, it sends an explicit request to obtain the full payload. This
design avoids sending the full message to every target immediately and
reduces unnecessary transmissions when the message has already reached a
node through another path.

The protocol also includes an adaptive mechanism based on duplicate
observations. Each node tracks how many first-time and duplicate full
messages it receives over a period of time. When the observed duplicate
rate is high, the eager fanout is reduced. When the duplicate rate is
low, the eager fanout is increased, within predefined bounds. The goal
is to adapt dissemination aggressiveness to the redundancy currently
observed in the network.

\subsection{Message Processing and Deduplication}

All broadcast protocols in our system use message identifiers to ensure
that each application message is delivered at most once per node. When a
message is received for the first time, it is stored in the local cache,
delivered to the application, and processed according to the semantics
of the selected broadcast protocol. If the same message is received
again, it is classified as a duplicate and discarded, although in the
hybrid protocol duplicate observations are also used as input for the
adaptive fanout mechanism.

This deduplication mechanism is fundamental for correctness. Without it,
both flood and gossip-based protocols would repeatedly retransmit the
same content, leading to unnecessary work and potentially unbounded
message circulation.

\subsection{Correctness Discussion}

At a high level, the correctness of the implementation depends on two
main properties. First, the membership layer must maintain a connected
enough overlay so that broadcast can progress through the network.
Second, the broadcast layer must guarantee that a message is delivered at
most once at each node while still being forwarded according to the
rules of the protocol.

The first property is supported by the design of Cyclon and HyParView,
which continuously refresh neighbor information and attempt to avoid
overlay fragmentation. The second property is guaranteed by the use of
unique message identifiers and local caches of already seen messages.

For the hybrid protocol, correctness additionally depends on the
interaction between lazy notifications and explicit requests. A lazy
notification never delivers a message by itself; it only informs a node
that the message exists. Delivery occurs only when the full payload is
received, either eagerly or as a reply to a request. This preserves the
same at-most-once delivery semantics as the other broadcast protocols.

\subsection{Implementation Summary}

In summary, the implementation was structured to support a clean
comparison between overlay construction and dissemination strategies.
Cyclon and HyParView provide two distinct membership mechanisms, while
flooding, eager-push gossip, and the proposed hybrid adaptive protocol
provide three dissemination strategies with different reliability and
overhead trade-offs. This modular structure also supports the
experimental methodology presented later in the paper.

\section{Experimental Evaluation}
\label{sec:evaluation}

\section{Conclusions}
\label{sec:conclusions}


\end{document}
\endinput
%%
%% End of file `sample-sigconf.tex'.
